% !TEX root = ../../tree_notes.tex
% =====================================================================
%  Chapter 7 -- Heaps and Priority Queues
%  Status: OUTLINE (not written yet). See structure/chapter.txt.
% =====================================================================
\bookchapter{Heaps and Priority Queues}{c07}
\chapterinfo{LC Medium -- CF 1600}{Outline}

\plannedtopic{The Binary Heap}{8}{%
  \item Complete binary tree stored in an array; children at $2i$, $2i+1$
  \item Min-heap and max-heap order property
  \item Sift-up and sift-down with diagrams of each swap
  \item Insert and extract-min in $O(\log n)$
}

\plannedtopic{Building and Using a Heap}{8}{%
  \item Floyd's build-heap in $O(n)$ (not $O(n\log n)$) and why
  \item \code{std::priority\_queue}, \code{push}, \code{pop}, \code{top}
  \item Custom comparators and \code{pair} keys
  \item Heap as a greedy engine
}

\plannedtopic{Heaps as Trees}{7}{%
  \item A heap is a complete binary tree, so its height is $\floor{\log_2 n}$
  \item Binomial and Fibonacci heaps (the idea, not the code)
  \item Leftist and skew heaps: mergeable heaps
  \item Why contests almost always use the plain binary heap
}

\plannedtopic{Classic Heap Problems}{9}{%
  \item $k$-th largest / smallest element
  \item Merge $k$ sorted lists
  \item Running median with two heaps
  \item Top-$k$ frequent elements
  \item Scheduling with a min-heap of finishing times
}

\closechapter
